and define the blocks
\begin{equation}\label{eq:blocks}
 \mathsf M=\sum_{i=0}^{4}x_i,
 \qquad
 \mathsf N=\sum_{i=5}^{8}x_i,
 \qquad
 \mathsf L=x_9+x_{10}.
\end{equation}

\classification{\newderivation}{Kernel checked in \path{Gate1B/R11/Polytope611.lean}.}
\begin{proposition}[Ordered-polytope geometry]\label{prop:611}
One has
\begin{equation}\label{eq:611-ledger}
 \mathsf M+\mathsf N+\mathsf L=1,
 \qquad \mathsf M\le\frac5{11},
 \qquad 1-\mathsf M\ge\frac6{11},
 \qquad \mathsf L\ge\frac2{11}.
\end{equation}
\end{proposition}

\begin{proof}
The block identity is the partition of the eleven coordinates.  Since every
one of the first five coordinates is at most every one of the remaining six,
summing the thirty pairwise inequalities gives
\[
            6\mathsf M\le5(1-\mathsf M),
\]
whence $11\mathsf M\le5$.  The complementary bound follows immediately.
Likewise, each of the first nine coordinates is at most each of the final two;
summing the eighteen inequalities gives
\[
            2(1-\mathsf L)\le9\mathsf L,
\]
so $\mathsf L\ge2/11$.
\end{proof}

The CARD5 chamber at parameter $\gamma$ consists of the two inequalities
\begin{equation}\label{eq:card5-chamber}
 x_6+x_7+x_8+x_9+x_{10}\le\gamma,
 \qquad
 x_0+x_7+x_8+x_9+x_{10}\le\gamma.
\end{equation}

\classification{\newderivation}{Kernel checked; the counterexamples are also kernel checked.}
\begin{proposition}[Pinned and unpinned chamber bounds]\label{prop:chamber}
Under \cref{eq:card5-chamber},
\begin{equation}\label{eq:chamber-unpinned}
 \mathsf L\le\gamma,
 \qquad
 \mathsf M\ge1-\frac65\gamma,
 \qquad
 \mathsf N\ge\frac23(1-\gamma).
\end{equation}
If in addition $\gamma\le\frac12-\eps$, then
\begin{equation}\label{eq:chamber-pinned}
 \mathsf M\ge\frac25+\frac65\eps,
 \qquad
 \mathsf N\ge\frac13+\frac23\eps.
\end{equation}
The two lower bounds in \cref{eq:chamber-pinned} do not follow from the chamber
conditions with a free $\gamma$.
\end{proposition}

\begin{proof}
The first chamber inequality gives $\mathsf L\le\gamma$.  Monotonicity gives
$5x_5\le x_6+\cdots+x_{10}\le\gamma$; hence
\[
 1-\mathsf M=x_5+\cdots+x_{10}\le x_5+\gamma
          \le\frac65\gamma.
\]
For the $\mathsf N$ bound, the first chamber inequality and total mass give
\[
       1-\gamma\le x_0+x_1+\cdots+x_5\le6x_5.
\]
Thus $x_5\ge(1-\gamma)/6$.  Since
$x_5\le x_6\le x_7\le x_8$, all four coordinates in $\mathsf N$ are at least
$x_5$, and hence
\[
       \mathsf N=x_5+x_6+x_7+x_8
          \ge4x_5\ge\frac23(1-\gamma).
\]
Substituting $\gamma\le1/2-\eps$ gives \cref{eq:chamber-pinned}.

For nonimplication, the formal bank provides two explicit ordered vectors.  In
the first, five coordinates equal $1/100$ and six equal $95/600$, so the chamber
holds at $\gamma=1$ but $\mathsf M=1/20<2/5$.  In the second, nine coordinates
equal $1/100$ and two equal $91/200$, so the chamber again holds at $\gamma=1$
but $\mathsf N=1/25<1/3$.
\end{proof}

\begin{remark}[Pascadi exponent dictionary]
The formal statement at $\theta=7/32$ is only finite rational arithmetic.  In
particular,
\[
                 2\theta=\frac7{16}<\frac6{11}\le1-\mathsf M.
\]
No analytic theorem of Pascadi is encoded or deduced from this inequality.  Any
later invocation of a Kloosterman theorem requires a separate coefficient,
modulus, and range dictionary \cite{Pascadi2025}.
\end{remark}

\section{Vaughan and long-M\"obius chart equivalence}\label{sec:vaughan}

Write $\zeta$ for the constant-one arithmetic function, and for an arithmetic
function $f$ let $f_{\le U}$ and $f_{>U}$ denote the exact complementary
truncations.  Vaughan's identity is used only as an algebraic decomposition;
its classical provenance may be found in \cite{Vaughan1997}.

\classification{\newreduction}{Kernel checked in \path{Gate1B/R11/VaughanCharts.lean}.}
\begin{theorem}[Exact four-lane identity and $V=2$ chart equivalence]\label{thm:vaughan}
In the convolution ring of real-valued arithmetic functions,
\begin{equation}\label{eq:vaughan}
 \Lambda
 =\Lambda_{\le V}+\mu_{\le U}*\log
  -\mu_{\le U}*\Lambda_{\le V}*\zeta
  +\mu_{>U}*\Lambda_{>V}*\zeta.
\end{equation}
For odd $N$ and $V=2$, the first and third lanes vanish, the second lane is the
low M\"obius--log block, and the fourth lane is the long M\"obius--log block.
Consequently,
\begin{equation}\label{eq:chart-equivalence}
       \operatorname{FourLaneValue}(U,2;N)
       =\operatorname{LongMobiusValue}(U;N)
       =\Lambda(N),
\end{equation}
with literal equality and no discarded mass.
\end{theorem}

\begin{proof}
Use $\mu*\zeta=\mathbf1$ and $\Lambda*\zeta=\log$.  Splitting
$\mu=\mu_{\le U}+\mu_{>U}$ and
$\Lambda=\Lambda_{\le V}+\Lambda_{>V}$ gives
\[
 \Lambda=(\mu_{\le U}+\mu_{>U})*\log.
\]
Insert $\log=\Lambda*\zeta$ in the high part and add and subtract
$\mu_{\le U}*\Lambda_{\le V}*\zeta$ to obtain \cref{eq:vaughan}.  When $N$ is
odd, every divisor of $N$ is odd.  The only odd integer at most $2$ is $1$, and
$\Lambda(1)=0$; hence $\Lambda_{\le2}$ vanishes on every divisor appearing in
the first and third lanes.  Moreover
$(\Lambda_{>2}*\zeta)(N)=\log N$ for odd $N$.  This identifies the remaining
two lanes with the low and long M\"obius--log sums.  The formal proof also
shows independently that both chart values equal $\Lambda(N)$.
\end{proof}

\classification{\newreduction}{Kernel checked in \path{Gate1B/R11/LongMobius.lean}.}
\begin{proposition}[Long-M\"obius divisor involution]\label{prop:long-reindex}
For $N\ne0$,
\begin{equation}\label{eq:long-reindex}
 \sum_{\substack{d\mid N\\d>U}}\mu(d)\log(N/d)
 =\sum_{\substack{k\mid N\\N/k>U}}\mu(N/k)\log k.
\end{equation}
The term $k=1$ vanishes exactly.  The same identity holds after multiplication
by an arbitrary finitely supported outer weight and summation over $N=n+2$.
\end{proposition}

\begin{proof}
The map $d\mapsto N/d$ is an involutive bijection between the two displayed
sets of divisors.  Substitution gives \cref{eq:long-reindex}; the $k=1$ term is
zero because $\log1=0$.
\end{proof}

\begin{warning}[What chart equivalence does not prove]
Neither \cref{thm:vaughan} nor \cref{prop:long-reindex} is an analytic
cancellation estimate.  They guarantee that changing charts loses no mass; the
long-M\"obius residual still requires a separate analytic argument.
\end{warning}

\section{The affine determinant line, reciprocity, and Bruhat structure}\label{sec:affine}

The matched source equation is
\begin{equation}\label{eq:matched}
                    AB+2=kd,
 \qquad\text{equivalently}\qquad AB-kd=-2.
\end{equation}
For odd $A$ and $B$, every cross-gcd is one:
\begin{equation}\label{eq:four-gcds}
 (A,k)=(A,d)=(B,k)=(B,d)=1.
\end{equation}
Indeed, any common divisor of a cross-pair divides $2$ and also an odd member
of the pair.

\classification{\newreduction}{Kernel checked in \texttt{Determinant.lean} and \texttt{AffineParametrization.lean}.}
\begin{theorem}[Unique affine parametrisation]\label{thm:affine}
Let $k\ne0$, $(A,k)=1$, and let $(B_0,d_0)$ be one integer solution of
$AB+2=kd$.  Then
\begin{equation}\label{eq:affine-biconditional}
 AB+2=kd
 \quad\Longleftrightarrow\quad
 \exists!\,t\in\Z:
 \begin{cases}
 B=B_0+kt,\\
 d=d_0+At.
 \end{cases}
\end{equation}
For two points separated by $h$ on this line,
\begin{equation}\label{eq:det2h}
             d_tB_{t+h}-d_{t+h}B_t=2h.
\end{equation}
\end{theorem}

\begin{proof}
Subtract the base equation from $AB+2=kd$ to obtain
\[
            A(B-B_0)=k(d-d_0).
\]
Since $(A,k)=1$, divisibility gives $k\mid B-B_0$, say $B-B_0=kt$.
Cancelling $k\ne0$ then gives $d-d_0=At$.  Uniqueness follows from
$k(t-s)=0$.  Conversely, substitution proves that every $t$ gives a solution.
Finally,
\begin{align*}
 d_tB_{t+h}-d_{t+h}B_t
 &=d_t(B_t+kh)-(d_t+Ah)B_t\\
 &=h(kd_t-AB_t)=2h,
\end{align*}
by \cref{eq:matched}.
\end{proof}

The scale relations \cref{eq:conservation-laws} are exact at exponent level.
The formal bank also proves the useful hybrid inequality: if $KT=\BB$ and
$a,b>0$, then
\begin{equation}\label{eq:hybrid-min}
             \min(T^{-a},K^{-b})\le\BB^{-ab/(a+b)}.
\end{equation}
It follows by separating the cases $T\ge\BB^{b/(a+b)}$ and its complement.

\classification{\newreduction}{Kernel checked in \path{Gate1B/R11/Reciprocity.lean}.}
\begin{proposition}[Exact reciprocity at shift two]\label{prop:reciprocity}
Let $(A,k)=1$, and choose $A',k'$ with
$AA'\equiv1\pmod{k}$ and $kk'\equiv1\pmod A$.  Then there is an integer
$n$ such that
\begin{equation}\label{eq:reciprocity-real}
       \frac{A'}k+\frac{k'}A=\frac1{Ak}+n,
\end{equation}
and for every integer $h$,
\begin{equation}\label{eq:reciprocity-exp}
       \eord{k}{-2hA'}
       =\eord{A}{2hk'}\eord{Ak}{-2h}.
\end{equation}
\end{proposition}

\begin{proof}
The congruences imply that both $k$ and $A$ divide
$AA'+kk'-1$.  Coprimality therefore gives
$Ak\mid AA'+kk'-1$, yielding \cref{eq:reciprocity-real}.  Multiply by
$-2h$ and exponentiate modulo one to obtain \cref{eq:reciprocity-exp}.
\end{proof}

\classification{\newreduction}{Kernel checked in \path{Gate1B/R11/Bruhat.lean}.}
\begin{proposition}[Exact $\mathrm{SL}_2$ Bruhat factorisation]\label{prop:bruhat}
Let
\[
 M=\begin{pmatrix}A&k\\d&B\end{pmatrix},\qquad AB-kd=-2,
\]
choose $\lambda$ with $\lambda(-2)=1$, and suppose $d^{-1}$ is supplied in the
coefficient ring.  Then
\begin{equation}\label{eq:bruhat}
 \begin{pmatrix}\lambda A&\lambda k\\d&B\end{pmatrix}
 =U(\lambda A d^{-1})D(d^{-1})S U(Bd^{-1}),
\end{equation}
